\section{Гидродинамика}

\subsection{Подходы Лагранжа и Эйлера}
Существует два основных способа описания движения жидкости:
\begin{itemize}
    \item \textbf{Метод Лагранжа:} рассматривает движение каждой отдельной частицы жидкости. Мы следим за траекторией, скоростью и ускорением конкретной частицы во времени.
    \item \textbf{Метод Эйлера:} фиксирует внимание на определенных точках пространства. Мы изучаем \textit{поле вектора скоростей} $\vec{v}(x, y, z, t)$, то есть какая скорость у жидкости в точке пространства в данный момент времени. 
\end{itemize}

\begin{center}
\begin{tblr}{
  colspec = {X[c]X[c]},
  vlines,
  hlines,
  row{1} = {font=\bfseries,mode=text},
  row{odd} = {bg=gray!10},
  cells = {c,m},
  rowsep = 6pt
}
\textbf{Подход Лагранжа} & \textbf{Подход Эйлера} \\
% --- Первый рисунок: Лагранж ---
\begin{tikzpicture}
    % Траектория
    \draw[thick, black, dashed, ->] (0.2,0.5) .. controls (1,1) and (1.5,2.5) .. (3,2.5);
    
    % Объем t1
    \draw[pattern=vertical lines, draw=black, thick] (0.2,0.5) circle (0.4);
    \node[fill=white, inner sep=1pt, scale=0.7] at (0.2, 0.5) {$t_1$};
    
    % Объем t2 (деформация)
    \draw[pattern=north east lines, draw=black, thick, rotate around={30:(1.5,1.8)}] (1.5,1.8) ellipse (0.5 and 0.3);
    \node[fill=white, inner sep=1pt, scale=0.7] at (1.5, 1.8) {$t_2$};
    
    % Объем t3
    \draw[pattern=dots, draw=black, thick, rotate around={-10:(3.2,2.5)}] (3.2,2.5) ellipse (0.7 and 0.2);
    \node[fill=white, inner sep=1pt, scale=0.7] at (3.2, 2.5) {$t_3$};
\end{tikzpicture} 
& 
% --- Второй рисунок: Эйлер ---
\begin{tikzpicture}
    % Линии тока
    \foreach \y in {0.5, 1.0, ..., 2.5} {
        \draw[black!40, -{Latex[length=1.5mm]}] (-0.5, \y) .. controls (1, \y+0.2) .. (3.5, \y-0.1);
    }
    
    % Контрольный объем
    \draw[pattern=crosshatch dots, draw=black, ultra thick, dashed] (1,0.8) rectangle (2,2.2);
    
    % Векторы скорости
    \draw[black, -Stealth, thick] (1.2,1.2) -- (1.8,1.3);
    \draw[black, -Stealth, thick] (1.2,1.8) -- (1.8,1.9);
    \draw[black, -Stealth, ultra thick] (1.5,1.5) -- (2.3,1.55) node[right, scale=0.7] {$\vec{v}(\vec{r}, t)$};
\end{tikzpicture} \\
Фиксируем поток частиц & Считаем поток частиц \\
Считаем координаты & Фиксируем координаты \\
\end{tblr}
\end{center}

\subsection{Уравнение неразрывности}



\subsubsection{Общий случай}
Выражает закон сохранения массы в каждой точке потока:
\[
\frac{\partial \rho}{\partial t} + \text{div}(\rho \vec{v}) = 0
\]
Где:
\begin{itemize}
    \item $\frac{\partial \rho}{\partial t}$ --- изменение плотности со временем в данной точке;
    \item $\text{div}(\rho \vec{v})$ --- дивергенция плотности потока, показывающая «растекание» вещества из точки.
\end{itemize}

\subsubsection{Для несжимаемой жидкости}
Если жидкость нельзя сжать (как воду), то $\frac{\partial \rho}{\partial t} = 0$ и $\rho$ выносится за знак дивергенции:
\[
\text{div} \vec{v} = 0
\]
Это означает, что поле вектора скоростей не имеет «источников» и «стоков»: сколько втекло в объем, столько и вытекло. Таким образом, поле скоростей несжимаемой жидкости является \textbf{соленоидальным}.

\subsubsection{Для трубки тока}
Для потока несжимаемой жидкости через трубу с меняющимся сечением $S$:
\[
S_1 v_1 = S_2 v_2 = \text{const}
\]
\textit{Вывод:} чем меньше сечение трубы, тем выше скорость потока в этом месте.

\subsubsection{Угловое колено трубы}

Рассмотрим угловое колено трубы, повёрнутое на $90^\circ$. 
Сверху в колено входит поток жидкости плотности \( \rho \) со скоростью $V_1$ через сечение площадью $S_1$, 
а выходит горизонтально со скоростью $V_2$ через сечение площадью $S_2$. 
Определить вектор силы $\vec{R}$ реакции опоры, действующей на колено:

\begin{center}
\begin{tikzpicture}
    % Труба с штриховкой
    \fill[pattern=north east lines] (0,2) -- (1,2) .. controls (1,1.3) and (1.3,1) .. (2,1) -- (2,0) .. controls (0,0) and (0,2) .. cycle;
    \draw[thick] (0,2) -- (1,2) .. controls (1,1.3) and (1.3,1) .. (2,1) -- (2,0) .. controls (0,0) and (0,2) .. cycle;
    
    % Скорости
    \draw[line width=2pt, -{Stealth[length=4pt,width=4pt]}] (0.5, 2.5) -- (0.5, 2.1) node[midway, above=6pt] {$\vec{v}_1, S_1$};
    \draw[line width=2pt, -{Stealth[length=4pt,width=4pt]}] (2.1, 0.5) -- (2.5, 0.5) node[midway, right=6pt] {$\vec{v}_2, S_2$};
    
    % Отражение сил (динамических)
    \draw[line width=1pt, dashed, -{Stealth[length=4pt,width=4pt]}] (1, 1) -- (2, 2) node[above right] {$\vec{F}$};
    \draw[line width=1pt, densely dotted, -{Stealth[length=4pt,width=4pt]}] (1, 1) -- (0, 0) node[below left] {$\vec{R}$};

    \node[circle, fill=black, inner sep=1.5pt] at (1,1) {};
    
    % Оси
    \draw[->, thin] (-0.5, 0) -- (3, 0) node[right] {$x$};
    \draw[->, thin] (0, -0.5) -- (0, 3) node[above] {$y$};
\end{tikzpicture}
\end{center}

Для несжимаемой жидкости массовый расход $\mu$ остается постоянным во всех сечениях:
\[
\mu = \frac{dm}{dt} = \rho S_1 v_1 = \rho S_2 v_2
\]

Согласно второму закону Ньютона в форме Мещерского для стационарного потока, сила $\vec{F}_{flow}$, действующая на жидкость со стороны стенок, равна изменению потока импульса:
\[
\vec{F} = \mu \vec{v}_2 - \mu \vec{v}_1
\]

Сила реакции $\vec{R}$, действующая на трубу со стороны жидкости, по третьему закону Ньютона равна:
\[
\vec{R} = -\vec{F} = \mu \vec{v}_1 - \mu \vec{v}_2 = \mu (\vec{v}_1 - \vec{v}_2)
\]

Подставляя выражение для массового расхода $\mu = \rho S_1 v_1$ и учитывая связь модулей скоростей через неразрывность ($v_2 = v_1 \frac{S_1}{S_2}$), запишем силу реакции через параметры входного потока:
\[
\boxed{\vec{R} = \rho S_1 v_1 \left( \vec{v}_1 - \vec{v}_2 \right)}
\]

\subsection{Поле вектора скоростей и плотность потока}

\textbf{Поле вектора скоростей} $\vec{v}(\vec{r}, t)$ --- распределение векторов скорости во всем объёме жидкости в момент времени \( t \).

\textbf{Стационарный поток} --- течение, при котором параметры жидкости (скорость $\vec{v}$, давление $P$, плотность $\rho$) в каждой точке пространства не изменяются с течением времени:
\[
\frac{\partial \vec{v}}{\partial t} = 0, \quad \frac{\partial P}{\partial t} = 0
\]
В таком потоке линии тока (траектории частиц) остаются неизменными во времени.

\textbf{Вектор плотности потока массы} $\vec{j}$ определяется как:
\[
\vec{j} = \rho \vec{v}
\]

Его направление совпадает с направлением скорости, а модуль равен массе, проходящей через единичную площадку в секунду. 

\subsection{Уравнение Бернулли}

\subsubsection{Формулировка}

\textbf{Невязкая жидкость} --- физическая модель жидкости, в которой внутренним трением и теплопроводностью пренебрегают.

Для стационарного потока идеальной \textit{(невязкой и несжимаемой)} жидкости выполняется закон сохранения энергии, записанный в виде \textit{уравнения Бернулли}: 
\[
P + \rho g h + \frac{\rho v^2}{2} = \text{const}
\]
Где:
\begin{itemize}
    \item $P$ --- статическое давление; 
    \item $\rho g h$ --- гидростатическое давление; 
    \item $\frac{\rho v^2}{2}$ --- динамическое давление. 
\end{itemize}

\subsubsection{Вывод}

Рассмотрим установившееся течение идеальной несжимаемой жидкости в трубке тока. Выделим малый объём жидкости, ограниченный сечениями $S_1$ и $S_2$. За время $\Delta t$ этот объём сместится: сечение $S_1$ переместится на $\Delta l_1 = v_1 \Delta t$, а сечение $S_2$ — на $\Delta l_2 = v_2 \Delta t$.

Так как течение стационарное, энергия средней части потока не изменяется со временем. Изменение \textit{полной механической энергии} выделенного объёма складывается из изменения кинетической и потенциальной энергии:

\[
\Delta E = \Delta E_k + \Delta E_p
\]

Масса малого элемента $\Delta m = \rho \Delta V$, где $\Delta V = S_1 \Delta l_1 = S_2 \Delta l_2$ — объём, прошедший через сечение (по несжимаемости). Тогда:
\[
\Delta E_k = \frac{\Delta m}{2} (v_2^2 - v_1^2) = \frac{\rho \Delta V}{2} (v_2^2 - v_1^2)
\]
\[
\Delta E_p = \Delta m g (h_2 - h_1) = \rho \Delta V g (h_2 - h_1)
\]

Суммируя:
\[
\Delta E = \rho \Delta V \left[ \frac{1}{2}(v_2^2 - v_1^2) + g(h_2 - h_1) \right]
\]

Работа внешних сил (сил давления) над выделенным объёмом:
\[
A = A_1 + A_2 = P_1 S_1 \Delta l_1 - P_2 S_2 \Delta l_2 = (P_1 - P_2) \Delta V
\]

Согласно закону сохранения энергии, изменение полной энергии равно работе внешних сил:
\[
\Delta E = A
\]

Подставляем выражения:
\[
\rho \Delta V \left[ \frac{1}{2}(v_2^2 - v_1^2) + g(h_2 - h_1) \right] = (P_1 - P_2) \Delta V
\]

Делим на $\Delta V$ и преобразуем выражение:
\[
P_1 + \rho g h_1 + \frac{\rho v_1^2}{2} = P_2 + \rho g h_2 + \frac{\rho v_2^2}{2}
\]

Поскольку сечения $S_1$ и $S_2$ выбраны произвольно, вдоль линии тока для идеальной несжимаемой жидкости выполняется:

\[
\boxed{P + \rho g h + \frac{\rho v^2}{2} = \mathrm{const}}
\]

\subsection{Эффект Вентури }

\subsubsection{Формулировка}

\textbf{Эффект Вентури} --- физическое явление, при котором: 
\begin{enumerate}
    \item Согласно уравнению неразрывности, при сужении трубы (уменьшении $S$) скорость потока $v$ возрастает. 
    \item Согласно уравнению Бернулли, при увеличении скорости $v$ давление $P$ в этом месте падает. 
\end{enumerate}
\textit{Вывод:} в узкой части трубы скорость жидкости выше, а статическое давление --- ниже.

\subsubsection{Расход жидкости}

По горизонтальной трубе переменного сечения течет вода. Площадь широкой части трубы равна $S_1$, а узкой --- $S_2$. Разность уровней воды в манометрических трубках составляет $\Delta h$. Найдите объемный расход воды $Q$:

\begin{center}
\begin{tikzpicture}[scale=1.2, >=stealth]
    % Контур трубы
    \draw[line width=1pt] (0,1) -- (1.5,1) -- (2,0.6) -- (3,0.6) -- (3.5,1) -- (5,1);
    \draw[line width=1pt] (0,-1) -- (1.5,-1) -- (2,-0.6) -- (3,-0.6) -- (3.5,-1) -- (5,-1);

    % Заполнение трубы (очень легкий узор для обозначения жидкости)
    \fill[pattern=crosshatch dots, pattern color=black!10] 
        (0,-1) -- (1.5,-1) -- (2,-0.6) -- (3,-0.6) -- (3.5,-1) -- (5,-1) -- 
        (5,1) -- (3.5,1) -- (2,0.6) -- (1.5,1) -- (0,1) -- cycle;
    
    % Манометрические трубки (более толстые)
    \draw[line width=0.8pt] (0.5,1) -- (0.5,2.5);
    \draw[line width=0.8pt] (0.8,1) -- (0.8,2.5);
    \draw[line width=0.8pt] (2.35,0.6) -- (2.35,2.5);
    \draw[line width=0.8pt] (2.65,0.6) -- (2.65,2.5);
    
    % Вода в манометрах (узор вместо цвета)
    \fill[pattern=horizontal lines, pattern color=black!30] (0.52,1) rectangle (0.78,2.2);
    \fill[pattern=horizontal lines, pattern color=black!30] (2.37,0.6) rectangle (2.63,1.5);
    \draw[line width=1pt, densely dashed] (0.52,2.2) -- (0.78,2.2);
    \draw[line width=1pt, densely dashed] (2.37,1.5) -- (2.63,1.5);
    
    % Линии уровней и дельта h (штриховые линии)
    \draw[dashed, line width=0.6pt] (0.8,2.2) -- (2.8,2.2);
    \draw[dashed, line width=0.6pt] (2.65,1.5) -- (2.8,1.5);
    \draw[<->, line width=0.6pt] (2.75,1.5) -- (2.75,2.2) node[midway, right] {$\Delta h$};
    
    % Сечения и скорости (разные типы стрелок)
    \draw[<->, line width=0.6pt] (0.2,-0.9) -- (0.2,0.9) node[midway, left] {$S_1$};
    \draw[<->, line width=0.6pt] (2.5,-0.5) -- (2.5,0.5) node[midway, left] {$S_2$};
    
    % Скорости - разные типы стрелок для различия
    \draw[->, line width=1.5pt] (0.3,0) -- (0.8,0) node[above] {$\vec{v}_1$};
    \draw[->, line width=1.5pt, dash pattern=on 4pt off 2pt] (2.6,0) -- (3.3,0) node[above] {$\vec{v}_2$};
    
    % Подписи к манометрам
    \node[above] at (0.65,2.5) {М1};
    \node[above] at (2.5,2.5) {М2};
\end{tikzpicture}
\end{center}

Согласно уравнению неразрывности для несжимаемой жидкости:
\[ Q = S_1 v_1 = S_2 v_2 \implies v_2 = v_1 \frac{S_1}{S_2} \]

Запишем уравнение Бернулли для горизонтальной трубки ($h_1 = h_2$):
\[ P_1 + \frac{\rho v_1^2}{2} = P_2 + \frac{\rho v_2^2}{2} \]

Разность давлений $P_1 - P_2$ выражается через разность высот в манометрах:
\[ P_1 - P_2 = \rho g \Delta h \]

Подставим выражение для $v_2$ в уравнение Бернулли:
\[ \rho g \Delta h = \frac{\rho}{2} (v_2^2 - v_1^2) = \frac{\rho v_1^2}{2} \left[ \left( \frac{S_1}{S_2} \right)^2 - 1 \right] \]

Выразим скорость $v_1$:
\[ v_1 = \sqrt{\frac{2g \Delta h}{(S_1/S_2)^2 - 1}} = S_2 \sqrt{\frac{2g \Delta h}{S_1^2 - S_2^2}} \]

Итоговая формула для расхода $Q$:
\[ \boxed{Q = S_1 S_2 \sqrt{\frac{2g \Delta h}{S_1^2 - S_2^2}}} \]